% Appendix A — Mathematical derivations
%
% Demonstrates: appendix chapter structure, lengthy multi-step
% derivations with align, cross-references to equations defined in
% earlier chapters, and \notag to suppress intermediate labels.

\chapter{Mathematical derivations}
\addthumb{Appendix}{\thechapter}{white}{thesisblue}
\clearpage

This appendix contains the proofs of two fundamental results from
Fourier analysis that underlie the signal-processing content of
\cref{ch1/ch1}.
Both theorems follow from the Fourier transform pair defined in
\cref{eq:ch1-fourier-fwd,eq:ch1-fourier-inv} in Chapter~1.

\section{The convolution theorem}

\textbf{Theorem.}
\textit{Let $f, g \in L^{1}(\mathbb{R})$.
Then}
\begin{equation}
  \mathcal{F}\{f * g\}(\xi) = \hat{f}(\xi)\cdot\hat{g}(\xi),
  \label{eq:app-convthm}
\end{equation}
\textit{i.e., the Fourier transform maps convolution
(defined in~\eqref{eq:ch1-convolution}) to pointwise multiplication.}

\medskip
\textbf{Proof.}
Starting from the definition of the Fourier transform
(\cref{eq:ch1-fourier-fwd}) and substituting the convolution
(\cref{eq:ch1-convolution}):

\begin{align}
  \mathcal{F}\{f * g\}(\xi)
    &= \int_{-\infty}^{\infty}
         (f * g)(t)\,e^{-2\pi i\xi t}\,dt
       \label{eq:app-conv-step1} \\
    &= \int_{-\infty}^{\infty}
         \left(
           \int_{-\infty}^{\infty} f(\tau)\,g(t-\tau)\,d\tau
         \right)
         e^{-2\pi i\xi t}\,dt.
       \notag
\end{align}

Since $f, g \in L^{1}$, Fubini's theorem permits interchanging
the order of integration:

\begin{align}
  &= \int_{-\infty}^{\infty} f(\tau)
       \left(
         \int_{-\infty}^{\infty} g(t-\tau)\,
         e^{-2\pi i\xi t}\,dt
       \right) d\tau.
     \label{eq:app-conv-fubini}
\end{align}

Substitute $u = t - \tau$ (so $t = u + \tau$ and $dt = du$) in the
inner integral:

\begin{align}
  &= \int_{-\infty}^{\infty} f(\tau)
       \left(
         \int_{-\infty}^{\infty} g(u)\,
         e^{-2\pi i\xi(u+\tau)}\,du
       \right) d\tau \notag \\
  &= \int_{-\infty}^{\infty} f(\tau)\,e^{-2\pi i\xi\tau}
       \underbrace{%
         \left(
           \int_{-\infty}^{\infty} g(u)\,
           e^{-2\pi i\xi u}\,du
         \right)%
       }_{=\,\hat{g}(\xi)}
     d\tau
     \label{eq:app-conv-factor} \\
  &= \hat{g}(\xi)
     \underbrace{%
       \int_{-\infty}^{\infty} f(\tau)\,
       e^{-2\pi i\xi\tau}\,d\tau%
     }_{=\,\hat{f}(\xi)}
     \notag \\
  &= \hat{f}(\xi)\cdot\hat{g}(\xi).
     \label{eq:app-conv-result}
\end{align}

This establishes~\eqref{eq:app-convthm}.\hfill$\square$

\bigskip

A direct consequence is that filtering in the time domain (a
convolution) is equivalent to multiplication in the frequency
domain, which is why frequency-domain filtering is computationally
attractive: an $\mathcal{O}(n^{2})$ time-domain convolution becomes
an $\mathcal{O}(n\log n)$ operation via the Fast Fourier Transform.

\section{Parseval's theorem}

\textbf{Theorem (Parseval / Rayleigh).}
\textit{Let $f \in L^{2}(\mathbb{R})$.
Then}
\begin{equation}
  \int_{-\infty}^{\infty} \lvert f(t) \rvert^{2}\,dt
  =
  \int_{-\infty}^{\infty} \lvert \hat{f}(\xi) \rvert^{2}\,d\xi.
  \label{eq:app-parseval}
\end{equation}
\textit{The Fourier transform is therefore an isometry on
$L^{2}(\mathbb{R})$: it preserves the total energy of a signal.}

\medskip
\textbf{Proof.}
Write $\lvert f(t)\rvert^{2} = f(t)\,\overline{f(t)}$ and express
the complex conjugate $\overline{f(t)}$ via the inverse Fourier
transform (\cref{eq:ch1-fourier-inv}):

\begin{align}
  \overline{f(t)}
    &= \overline{
         \int_{-\infty}^{\infty}
           \hat{f}(\xi)\,e^{+2\pi i\xi t}\,d\xi
       }
     = \int_{-\infty}^{\infty}
         \overline{\hat{f}(\xi)}\,
         e^{-2\pi i\xi t}\,d\xi.
     \label{eq:app-conj}
\end{align}

Substituting~\eqref{eq:app-conj}:

\begin{align}
  \int_{-\infty}^{\infty} \lvert f(t)\rvert^{2}\,dt
    &= \int_{-\infty}^{\infty} f(t)
         \left(
           \int_{-\infty}^{\infty}
             \overline{\hat{f}(\xi)}\,
             e^{-2\pi i\xi t}\,d\xi
         \right) dt.
       \label{eq:app-pars-step1}
\end{align}

Interchanging the order of integration (Fubini's theorem, valid
since $f \in L^{2}$ and $\hat{f} \in L^{2}$):

\begin{align}
  &= \int_{-\infty}^{\infty} \overline{\hat{f}(\xi)}
       \underbrace{%
         \left(
           \int_{-\infty}^{\infty} f(t)\,
           e^{-2\pi i\xi t}\,dt
         \right)%
       }_{=\,\hat{f}(\xi)}
     d\xi
     \label{eq:app-pars-fubini} \\
  &= \int_{-\infty}^{\infty}
       \overline{\hat{f}(\xi)}\,\hat{f}(\xi)\,d\xi \notag \\
  &= \int_{-\infty}^{\infty}
       \lvert\hat{f}(\xi)\rvert^{2}\,d\xi,
     \label{eq:app-pars-result}
\end{align}

where the last step uses $\bar{z}z = \lvert z \rvert^{2}$ for any
$z \in \mathbb{C}$.
This establishes~\eqref{eq:app-parseval}.\hfill$\square$

\bigskip

Parseval's theorem has an important physical interpretation: the
total energy computed in the time domain equals the total energy
computed in the frequency domain.
In the context of the \ac{snr} definition used in Chapter~1,
this means that filtering operations that attenuate
noise in the frequency domain produce a provably equivalent reduction
in time-domain noise energy.
